\documentclass[10pt,a4paper]{amsart}
\usepackage{amsmath,amssymb,amsthm,graphicx}
% For commutative diagrams we use Xy-pic
\usepackage[all]{xy}

% We use 2cell for 2-commutative diagrams.
\xyoption{2cell}
% Theorem environments.
%
\theoremstyle{plain}
\newtheorem{theorem}[subsection]{Theorem}
\newtheorem{proposition}[subsection]{Proposition}
\newtheorem{lemma}[subsection]{Lemma}

\theoremstyle{definition}
\newtheorem{definition}[subsection]{Definition}
\newtheorem{example}[subsection]{Example}
\newtheorem{exercise}[subsection]{Exercise}
\newtheorem{situation}[subsection]{Situation}

\theoremstyle{remark}
\newtheorem{remark}[subsection]{Remark}
\newtheorem{remarks}[subsection]{Remarks}

\numberwithin{equation}{subsection}

% Macros
%
\def\lim{\mathop{\mathrm{lim}}\nolimits}
\def\colim{\mathop{\mathrm{colim}}\nolimits}
\def\Spec{\mathop{\mathrm{Spec}}}
\def\Hom{\mathop{\mathrm{Hom}}\nolimits}
\def\Ext{\mathop{\mathrm{Ext}}\nolimits}
\def\SheafHom{\mathop{\mathcal{H}\!\mathit{om}}\nolimits}
\def\SheafExt{\mathop{\mathcal{E}\!\mathit{xt}}\nolimits}
\def\Sch{\mathit{Sch}}
\def\Mor{\mathop{\mathrm{Mor}}\nolimits}
\def\Ob{\mathop{\mathrm{Ob}}\nolimits}
\def\Sh{\mathop{\mathit{Sh}}\nolimits}
\def\NL{\mathop{N\!L}\nolimits}
\def\CH{\mathop{\mathrm{CH}}\nolimits}
\def\proetale{{pro\text{-}\acute{e}tale}}
\def\etale{{\acute{e}tale}}
\def\QCoh{\mathit{QCoh}}
\def\Ker{\mathop{\mathrm{Ker}}}
\def\Im{\mathop{\mathrm{Im}}}
\def\Coker{\mathop{\mathrm{Coker}}}
\def\Coim{\mathop{\mathrm{Coim}}}

% Boxtimes
%
\DeclareMathSymbol{\boxtimes}{\mathbin}{AMSa}{"02}

%
% Macros for moduli stacks/spaces
%
\def\QCohstack{\mathcal{QC}\!\mathit{oh}}
\def\Cohstack{\mathcal{C}\!\mathit{oh}}
\def\Spacesstack{\mathcal{S}\!\mathit{paces}}
\def\Quotfunctor{\mathrm{Quot}}
\def\Hilbfunctor{\mathrm{Hilb}}
\def\Curvesstack{\mathcal{C}\!\mathit{urves}}
\def\Polarizedstack{\mathcal{P}\!\mathit{olarized}}
\def\Complexesstack{\mathcal{C}\!\mathit{omplexes}}
% \Pic is the operator that assigns to X its picard group, usage \Pic(X)
% \Picardstack_{X/B} denotes the Picard stack of X over B
% \Picardfunctor_{X/B} denotes the Picard functor of X over B
\def\Pic{\mathop{\mathrm{Pic}}\nolimits}
\def\Picardstack{\mathcal{P}\!\mathit{ic}}
\def\Picardfunctor{\mathrm{Pic}}
\def\Deformationcategory{\mathcal{D}\!\mathit{ef}}
\usepackage{cleveref}
\crefname{originallemma}{original lemma}{original lemmas}
\Crefname{originallemma}{Original lemma}{Original lemmas}
\makeatletter
\newlabel{lemma-compare-fppf-etale}{{075R}{source}{}{sourcepointer-1}{}}
\newlabel{lemma-compare-fppf-etale@cref}{{[originallemma][0][]075R}{[1][0][]source}{}{}{}}
\makeatother
\crefname{originallemma}{original lemma}{original lemmas}
\Crefname{originallemma}{Original lemma}{Original lemmas}
\makeatletter
\newlabel{lemma-compare-higher-direct-image-proper}{{0DDJ}{source}{}{sourcepointer-2}{}}
\newlabel{lemma-compare-higher-direct-image-proper@cref}{{[originallemma][0][]0DDJ}{[1][0][]source}{}{}{}}
\makeatother
\crefname{originallemma}{original lemma}{original lemmas}
\Crefname{originallemma}{Original lemma}{Original lemmas}
\makeatletter
\newlabel{lemma-push-pull-fppf-etale}{{0DDM}{source}{}{sourcepointer-3}{}}
\newlabel{lemma-push-pull-fppf-etale@cref}{{[originallemma][0][]0DDM}{[1][0][]source}{}{}{}}
\makeatother
\crefname{originallemma}{original lemma}{original lemmas}
\Crefname{originallemma}{Original lemma}{Original lemmas}
\makeatletter
\newlabel{sites-cohomology-lemma-A}{{0EZ9}{source}{}{sourcepointer-4}{}}
\newlabel{sites-cohomology-lemma-A@cref}{{[originallemma][0][]0EZ9}{[1][0][]source}{}{}{}}
\makeatother
\crefname{originallemma}{original lemma}{original lemmas}
\Crefname{originallemma}{Original lemma}{Original lemmas}
\makeatletter
\newlabel{sites-cohomology-lemma-V-C-all-n-general}{{0EZF}{source}{}{sourcepointer-5}{}}
\newlabel{sites-cohomology-lemma-V-C-all-n-general@cref}{{[originallemma][0][]0EZF}{[1][0][]source}{}{}{}}
\makeatother
\crefname{originallemma}{original lemma}{original lemmas}
\Crefname{originallemma}{Original lemma}{Original lemmas}
\makeatletter
\newlabel{sites-cohomology-lemma-V-implies-C-general}{{0EZA}{source}{}{sourcepointer-6}{}}
\newlabel{sites-cohomology-lemma-V-implies-C-general@cref}{{[originallemma][0][]0EZA}{[1][0][]source}{}{}{}}
\makeatother
\crefname{originallemma}{original lemma}{original lemmas}
\Crefname{originallemma}{Original lemma}{Original lemmas}
\makeatletter
\newlabel{sites-cohomology-lemma-derived-pushforward-composition}{{0D6E}{source}{}{sourcepointer-7}{}}
\newlabel{sites-cohomology-lemma-derived-pushforward-composition@cref}{{[originallemma][0][]0D6E}{[1][0][]source}{}{}{}}
\makeatother
\crefname{originalsection}{original section}{original sections}
\Crefname{originalsection}{Original section}{Original sections}
\makeatletter
\newlabel{sites-cohomology-section-compare-general}{{0EZ1}{source}{}{sourcepointer-8}{}}
\newlabel{sites-cohomology-section-compare-general@cref}{{[originalsection][0][]0EZ1}{[1][0][]source}{}{}{}}
\makeatother
\title{The V and C conditions for the comparison of the etale and fppf topologies, with the five comparison statements they yield}
\author{The Stacks Project Authors}
\date{stacks-project:etale-cohomology.tex:local-snapshot-20261011; GNU FDL 1.2 or later}
\begin{document}
\maketitle
\noindent\textit{Editorial note.} All delivered mathematical text is the author's own, copied from the pinned Stacks Project chapter file etale-cohomology.tex; no statement or proof step was rewritten, reordered or removed.\par\medskip
\noindent\textit{Editorial note.} The counted proof is the author's complete argument: the proof is the author's complete argument: each of the five parts is reduced to the corresponding general statement, and part (5) is carried out by the four-term chain of equalities plus the explicit induction on i showing that R epsilon\_\{Y,*\} reflects isomorphisms for maps out of a\_Y\textasciicircum{}\{-1\}L.\par\medskip
\noindent\textit{Editorial note.} The label lemma-V-C-all-n-etale-fppf names the counted statement, and the proof environment immediately following it in the source is delivered verbatim as its proof.\par\medskip
\noindent\textit{Editorial note.} The statement and proof are copied from the pinned chapter file; the eight cited inputs are kept as external pointers with their canonical Stacks tags.\par\medskip
\noindent\textit{Editorial note.} The preamble is the chapter's own native preamble excerpt already used for this source: the Xy-pic setup (native preamble lines 14-18) and the native theorem declarations with the mathematical macros (native preamble lines 154-217).\par\medskip
\noindent\textit{Allowed standard prerequisites (editorial).}\begin{itemize}
\item Comparison morphisms a\_X, pi\_X, epsilon\_X between the etale and fppf topoi over S. Same chapter, the morphisms of topoi relating X\_etale, (Sch/X)\_fppf and (Sch/X)\_etale and the composition a\_X = epsilon\_X composed with the comparison. Supplies every functor appearing in the five assertions.
\item The general V and C conditions of the sites chapter. Cohomology on Sites, the general lemmas proving the V-C conditions, the V-implies-C statement and the category A\_X of the comparison section. Supplies the abstract statements that the proof specialises to the etale-fppf setting.
\item Proper base change for torsion coefficients and comparison of higher direct images. Same chapter, the proper base change theorem and the comparison of the small and big higher direct images for proper morphisms. Supplies the third and fourth equalities in the chain for part (5).
\end{itemize}\medskip
\ifdefined\hypertarget\hypertarget{sourcepointer-1}{}\fi
\noindent\textit{Source reference.} Original lemma 075R: sources/stacks-etale-cohomology/pinned-original-source.tex, source lines 21666--21666. Prerequisite pointer invoked at proof line(s) 21749; it names a standard fact of the same chapter that the counted argument uses, and every step of the counted proof built on it is retained verbatim in the delivered body.\par\medskip
\ifdefined\hypertarget\hypertarget{sourcepointer-2}{}\fi
\noindent\textit{Source reference.} Original lemma 0DDJ: sources/stacks-etale-cohomology/pinned-original-source.tex, source lines 21323--21323. Prerequisite pointer invoked at proof line(s) 21794; it names a standard fact of the same chapter that the counted argument uses, and every step of the counted proof built on it is retained verbatim in the delivered body.\par\medskip
\ifdefined\hypertarget\hypertarget{sourcepointer-3}{}\fi
\noindent\textit{Source reference.} Original lemma 0DDM: sources/stacks-etale-cohomology/pinned-original-source.tex, source lines 21484--21484. Prerequisite pointer invoked at proof line(s) 21790; it names a standard fact of the same chapter that the counted argument uses, and every step of the counted proof built on it is retained verbatim in the delivered body.\par\medskip
\ifdefined\hypertarget\hypertarget{sourcepointer-4}{}\fi
\noindent\textit{Source reference.} Original lemma 0EZ9: sources/stacks-chow/tags.txt, source lines 18719--18719. Prerequisite pointer invoked at proof line(s) 21753; it names a standard fact of a sibling chapter that the counted argument uses, and every step of the counted proof built on it is retained verbatim in the delivered body.\par\medskip
\ifdefined\hypertarget\hypertarget{sourcepointer-5}{}\fi
\noindent\textit{Source reference.} Original lemma 0EZF: sources/stacks-chow/tags.txt, source lines 18725--18725. Prerequisite pointer invoked at proof line(s) 21759; it names a standard fact of a sibling chapter that the counted argument uses, and every step of the counted proof built on it is retained verbatim in the delivered body.\par\medskip
\ifdefined\hypertarget\hypertarget{sourcepointer-6}{}\fi
\noindent\textit{Source reference.} Original lemma 0EZA: sources/stacks-chow/tags.txt, source lines 18720--18720. Prerequisite pointer invoked at proof line(s) 21763; it names a standard fact of a sibling chapter that the counted argument uses, and every step of the counted proof built on it is retained verbatim in the delivered body.\par\medskip
\ifdefined\hypertarget\hypertarget{sourcepointer-7}{}\fi
\noindent\textit{Source reference.} Original lemma 0D6E: sources/stacks-chow/tags.txt, source lines 16519--16519. Prerequisite pointer invoked at proof line(s) 21792; it names a standard fact of a sibling chapter that the counted argument uses, and every step of the counted proof built on it is retained verbatim in the delivered body.\par\medskip
\ifdefined\hypertarget\hypertarget{sourcepointer-8}{}\fi
\noindent\textit{Source reference.} Original section 0EZ1: sources/stacks-chow/tags.txt, source lines 18711--18711. Prerequisite pointer invoked at proof line(s) 21750; it names a standard fact of a sibling chapter that the counted argument uses, and every step of the counted proof built on it is retained verbatim in the delivered body.\par\medskip
\begin{lemma}
\label{lemma-V-C-all-n-etale-fppf}
With notation as above.
\begin{enumerate}
\item For $X \in \Ob((\Sch/S)_{fppf})$ and an abelian sheaf $\mathcal{F}$
on $X_\etale$ we have
$\epsilon_{X, *}a_X^{-1}\mathcal{F} = \pi_X^{-1}\mathcal{F}$
and $R^i\epsilon_{X, *}(a_X^{-1}\mathcal{F}) = 0$ for $i > 0$.
\item For a finite morphism $f : X \to Y$ in $(\Sch/S)_{fppf}$
and abelian sheaf $\mathcal{F}$ on $X$ we have
$a_Y^{-1}(R^if_{small, *}\mathcal{F}) =
R^if_{big, fppf, *}(a_X^{-1}\mathcal{F})$
for all $i$.
\item For a scheme $X$ and $K$ in $D^+(X_\etale)$ the map
$\pi_X^{-1}K \to R\epsilon_{X, *}(a_X^{-1}K)$ is an isomorphism.
\item For a finite morphism $f : X \to Y$ of schemes
and $K$ in $D^+(X_\etale)$ we have
$a_Y^{-1}(Rf_{small, *}K) = Rf_{big, fppf, *}(a_X^{-1}K)$.
\item For a proper morphism $f : X \to Y$ of schemes
and $K$ in $D^+(X_\etale)$ with torsion cohomology sheaves we have
$a_Y^{-1}(Rf_{small, *}K) = Rf_{big, fppf, *}(a_X^{-1}K)$.
\end{enumerate}
\end{lemma}

\begin{proof}
By Lemma \ref{lemma-compare-fppf-etale} the lemmas in
Cohomology on Sites, Section \ref{sites-cohomology-section-compare-general}
all apply to our current setting. To translate the results
observe that the category $\mathcal{A}_X$ of
Cohomology on Sites, Lemma \ref{sites-cohomology-lemma-A}
is the essential image of
$a_X^{-1} : \textit{Ab}(X_\etale) \to \textit{Ab}((\Sch/X)_{fppf})$.

\medskip\noindent
Part (1) is equivalent to $(V_n)$ for all $n$ which holds by
Cohomology on Sites, Lemma \ref{sites-cohomology-lemma-V-C-all-n-general}.

\medskip\noindent
Part (2) follows by applying $\epsilon_Y^{-1}$ to the conclusion of
Cohomology on Sites, Lemma \ref{sites-cohomology-lemma-V-implies-C-general}.

\medskip\noindent
Part (3) follows from Cohomology on Sites, Lemma
\ref{sites-cohomology-lemma-V-C-all-n-general} part (1)
because $\pi_X^{-1}K$ is in $D^+_{\mathcal{A}'_X}((\Sch/X)_\etale)$
and $a_X^{-1} = \epsilon_X^{-1} \circ a_X^{-1}$.

\medskip\noindent
Part (4) follows from Cohomology on Sites, Lemma
\ref{sites-cohomology-lemma-V-C-all-n-general} part (2)
for the same reason.

\medskip\noindent
Part (5). We use that
\begin{align*}
R\epsilon_{Y, *}Rf_{big, fppf, *}a_X^{-1}K
& =
Rf_{big, \etale, *}R\epsilon_{X, *}a_X^{-1}K \\
& =
Rf_{big, \etale, *}\pi_X^{-1}K \\
& =
\pi_Y^{-1}Rf_{small, *}K \\
& =
R\epsilon_{Y, *} a_Y^{-1}Rf_{small, *}K
\end{align*}
The first equality by the commutative diagram in
Lemma \ref{lemma-push-pull-fppf-etale}
and Cohomology on Sites, Lemma
\ref{sites-cohomology-lemma-derived-pushforward-composition}.
The second equality is (3). The third is
Lemma \ref{lemma-compare-higher-direct-image-proper} part (2).
The fourth is (3) again. Thus the base change map
$a_Y^{-1}(Rf_{small, *}K) \to Rf_{big, fppf, *}(a_X^{-1}K)$
induces an isomorphism
$$
R\epsilon_{Y, *}a_Y^{-1}Rf_{small, *}K \to
R\epsilon_{Y, *}Rf_{big, fppf, *}a_X^{-1}K
$$
The proof is finished by the following remark: a map
$\alpha : a_Y^{-1}L \to M$ with $L$ in $D^+(Y_\etale)$
and $M$ in $D^+((\Sch/Y)_{fppf})$ such that $R\epsilon_{Y, *}\alpha$
is an isomorphism, is an isomorphism. Namely, 
we show by induction on $i$ that $H^i(\alpha)$ is an isomorphism.
This is true for all sufficiently small $i$.
If it holds for $i \leq i_0$, then we see that
$R^j\epsilon_{Y, *}H^i(M) = 0$ for $j > 0$ and $i \leq i_0$
by (1) because $H^i(M) = a_Y^{-1}H^i(L)$ in this range.
Hence $\epsilon_{Y, *}H^{i_0 + 1}(M) = H^{i_0 + 1}(R\epsilon_{Y, *}M)$
by a spectral sequence argument.
Thus $\epsilon_{Y, *}H^{i_0 + 1}(M) = \pi_Y^{-1}H^{i_0 + 1}(L) =
\epsilon_{Y, *}a_Y^{-1}H^{i_0 + 1}(L)$.
This implies $H^{i_0 + 1}(\alpha)$ is an isomorphism
(because $\epsilon_{Y, *}$ reflects isomorphisms as it is the
identity on underlying presheaves) as desired.
\end{proof}

\end{document}
